% =================================================
% Figures used across the "Background" section slides.
% Requires: tikz (with arrows.meta, calc, decorations.pathmorphing),
% xcolor. Included from main.tex via % =================================================
% Figures used across the "Background" section slides.
% Requires: tikz (with arrows.meta, calc, decorations.pathmorphing),
% xcolor. Included from main.tex via % =================================================
% Figures used across the "Background" section slides.
% Requires: tikz (with arrows.meta, calc, decorations.pathmorphing),
% xcolor. Included from main.tex via % =================================================
% Figures used across the "Background" section slides.
% Requires: tikz (with arrows.meta, calc, decorations.pathmorphing),
% xcolor. Included from main.tex via \input{figures.tex}.
% =================================================

% -----------------------------------------------
% Colors for feature-space / hypothesis-space sketches
% -----------------------------------------------
% class0/class1 = the two classes (o / x) everywhere below; hlA/hlB = the
% two highlight-ring colors used for x^(1)/x^(2) on the Statistical Model
% slide; testred = test points, decision-boundary curves, and histogram
% bars. Change these once here to re-theme every plot in the deck.
\definecolor{class0}{RGB}{56,142,60}    % green: y = 0, "o"
\definecolor{class1}{RGB}{81,45,168}    % purple: y = 1, "x"
\definecolor{hlA}{RGB}{216,27,96}       % pink highlight
\definecolor{hlB}{RGB}{230,81,0}        % orange highlight
\definecolor{testred}{RGB}{198,40,40}   % red: test point / boundary

% -----------------------------------------------
% Reusable TikZ pieces
% -----------------------------------------------
% Background point cloud shared by the two feature-space plots (Statistical
% Model / SL Framework frames): a class-0 ("o") cluster top-right, a
% class-1 ("x") cluster bottom-left.
%
% HOW TO ADJUST:
%   - The two \fill ellipses are just the soft shaded "cloud" background;
%     their (center) and (radius and radius) control where each cluster
%     sits and how big/spread-out it looks. They don't need to track the
%     individual points below -- they're a backdrop, not a fit.
%   - Each \node ... {$\circ$} / {$\times$} is ONE data point: the
%     coordinate after "at" is its (x2,x1) position (the plot has x2
%     horizontal, x1 vertical -- see the axis \draw's below). Add, remove,
%     or move these to change what D_n "looks like".
%   - The two \draw[-{Latex}] lines are the axes: change the second
%     coordinate to lengthen/shorten an axis; change the node[right]/
%     [above] text to relabel it.
%   - If you move anything outside roughly x2,x1 in [-0.5,6.4] x
%     [-0.5,4.6], also widen \fixedbbox below, or the picture will
%     silently resize and the plot will start "jumping" between slides
%     again (see \fixedbbox's own comment).
\newcommand{\featurecloud}{
  \fill[class0!20] (2.6,2.7) ellipse (1.35 and 1.05);
  \fill[class1!20] (1.0,1.0) ellipse (1.35 and 1.05);
  \node[class0] at (2.2,3.1) {$\circ$};
  \node[class0] at (2.6,3.35) {$\circ$};
  \node[class0] at (3.05,3.05) {$\circ$};
  \node[class0] at (2.35,2.65) {$\circ$};
  \node[class0] at (2.95,2.5) {$\circ$};
  \node[class0] at (3.35,2.85) {$\circ$};
  \node[class1] at (0.5,1.4) {$\times$};
  \node[class1] at (0.8,0.8) {$\times$};
  \node[class1] at (1.15,1.6) {$\times$};
  \node[class1] at (0.4,0.6) {$\times$};
  \node[class1] at (1.3,1.0) {$\times$};
  \node[class1] at (0.95,1.9) {$\times$};
  \draw[-{Latex}] (-0.3,0) -- (4.3,0) node[right] {$x_2$};
  \draw[-{Latex}] (0,-0.3) -- (0,4.3) node[above] {$x_1$};
}

% Small "$y=0$ / $y=1$" key box placed next to \featurecloud.
% HOW TO ADJUST:
%   - "5.1" in "(5.1+#1,1.7)" is the box's horizontal position (push it
%     further right if you widen \featurecloud, e.g. add more points out
%     past x2=4.3); "1.7" is its vertical position; #1 is an optional
%     extra offset a caller can pass, e.g. \featurelegend[0.3], without
%     editing this macro.
%   - "inner sep=4pt" / "font=\footnotesize" control the box's padding
%     and text size.
\newcommand{\featurelegend}[1][0]{
  \node[draw, align=left, font=\footnotesize, fill=white, inner sep=4pt]
    at (5.1+#1,1.7) {{\color{class0}$\circ$}\; $y=0$ \\ {\color{class1}$\times$}\; $y=1$};
}

% Invisible path that pins the bounding box of the feature-space plots to a
% fixed rectangle, so the picture is the same size on every slide that
% uses it (Statistical Model / SL Framework) regardless of which extra
% annotations (highlight circles, the test star, ...) are drawn on top --
% this is what keeps the plot from visually jumping when you flip between
% those slides.
%
% HOW TO ADJUST:
%   - Arguments are (xmin,ymin) rectangle (xmax,ymax), same (x2,x1)
%     coordinates as \featurecloud. Call this FIRST inside every
%     tikzpicture that uses \featurecloud, so all of them agree on one
%     size.
%   - If a new annotation you add (a star, a label, ...) sits outside
%     this rectangle, the picture's natural bounding box will grow to fit
%     it on THAT slide only, breaking the "same position everywhere"
%     guarantee. Widen the rectangle here (once, for every slide) instead
%     of nudging it per-frame.
\newcommand{\fixedbbox}{\path (-0.5,-0.5) rectangle (6.4,4.6);}

% "Dataset + decision boundary" icons used on the hypothesis-space slides
% (Meta Level bar charts, and the inline $\varphi_1,\varphi_2$ icons on
% the PPD slide): each hypothesis $\varphi$ is a *different* data
% generation mechanism, so each icon below has its own point-cloud
% geometry; the red curve just traces the shape of that particular cloud
% rather than being an independent line drawn over one fixed background.
%
% HOW TO ADJUST (applies to \hypIconBase and all four icons below it):
%   - All coordinates live in a fixed 1.4 x 1.4 box (see the \clip and
%     the border \draw here) -- keep new content roughly inside
%     [0,1.4] x [0,1.4] or it gets clipped off.
%   - #1 is the icon-specific content (fills, points, boundary curve),
%     supplied by each \hyp... macro below.
%   - These icons are always used heavily scaled down (via \scalebox or
%     inside \metapanel), so keep shapes simple/high-contrast -- fine
%     detail won't survive the shrink.
\newcommand{\hypIconBase}[1]{%
  \begin{tikzpicture}[baseline=(current bounding box.center)]
    \clip (0,0) rectangle (1.4,1.4);
    #1
    \draw (0,0) rectangle (1.4,1.4);
  \end{tikzpicture}%
}

% Main-diagonal split: "o" cluster top-right, "x" cluster bottom-left.
% HOW TO ADJUST:
%   - The two \fill ellipses are this hypothesis's cluster "shape"
%     (center + radii) -- move/resize them to change what the data looks
%     like under this $\varphi$.
%   - Each \node {$\circ$}/{$\times$} is one visible point; the final
%     \draw is the decision-boundary line -- its endpoints should track
%     wherever you put the two clusters (i.e. redraw it if you move the
%     ellipses, so it still separates them).
\newcommand{\hypLinA}{\hypIconBase{
  % half-space fills (clipped to box by \hypIconBase)
  \fill[class0!25] (-0.15,1.25) -- (1.55,0.15) -- (1.55,1.55) -- (-0.15,1.55) -- cycle;
  \fill[class1!25] (-0.15,1.25) -- (1.55,0.15) -- (1.55,-0.1) -- (-0.15,-0.1) -- cycle;
  \node[class0, font=\tiny] at (0.8,1.15) {$\circ$};
  \node[class0, font=\tiny] at (1.15,1.05) {$\circ$};
  \node[class0, font=\tiny] at (0.95,0.85) {$\circ$};
  \node[class1, font=\tiny] at (0.3,0.5) {$\times$};
  \node[class1, font=\tiny] at (0.6,0.3) {$\times$};
  \node[class1, font=\tiny] at (0.35,0.25) {$\times$};
  \draw[testred, thick] (-0.15,1.25) -- (1.55,0.15);
}}

% Anti-diagonal split: "x" cluster top-left, "o" cluster bottom-right --
% a genuinely different data layout from \hypLinA, not just a new line.
% HOW TO ADJUST: same recipe as \hypLinA (ellipses = clusters, nodes =
% points, final \draw = boundary line through the gap between them).
\newcommand{\hypLinB}{\hypIconBase{
  % half-space fills (clipped to box by \hypIconBase)
  \fill[class1!25] (-0.15,0.15) -- (1.55,1.25) -- (1.55,1.55) -- (-0.15,1.55) -- cycle;
  \fill[class0!25] (-0.15,0.15) -- (1.55,1.25) -- (1.55,-0.1) -- (-0.15,-0.1) -- cycle;
  \node[class1, font=\tiny] at (0.3,1.15) {$\times$};
  \node[class1, font=\tiny] at (0.6,1.05) {$\times$};
  \node[class1, font=\tiny] at (0.4,0.85) {$\times$};
  \node[class0, font=\tiny] at (1.1,0.55) {$\circ$};
  \node[class0, font=\tiny] at (0.85,0.25) {$\circ$};
  \node[class0, font=\tiny] at (1.15,0.2) {$\circ$};
  \draw[testred, thick] (-0.15,0.15) -- (1.55,1.25);
}}

% Quadratic split (y ~ x^2): × fills the bottom valley, ○ fills above.
% Boundary is the exact parabola y = 2(x-0.7)^2 + 0.35 (vertex at (0.7,0.35),
% arms reach y≈1.33 at the sides of the 1.4×1.4 box).
% HOW TO ADJUST:
%   - Vertex position: change (0.7) and (0.35) in the formula.
%   - Curvature: change the leading coefficient (currently 2).
%   - Data points: × must satisfy y < 2(x-0.7)^2+0.35; ○ must be above.
\newcommand{\hypCurve}{\hypIconBase{
  % half-space fills using the exact parabola (clipped to box by \hypIconBase)
  \fill[class0!25] plot[domain=-0.15:1.55, samples=60] (\x, {2*(\x-0.7)*(\x-0.7)+0.35})
    -- (1.55,1.55) -- (-0.15,1.55) -- cycle;
  \fill[class1!25] plot[domain=-0.15:1.55, samples=60] (\x, {2*(\x-0.7)*(\x-0.7)+0.35})
    -- (1.55,-0.1) -- (-0.15,-0.1) -- cycle;
  % × (class 1) -- below the parabola, spread across bottom valley
  \node[class1, font=\tiny] at (0.7,0.15)  {$\times$};
  \node[class1, font=\tiny] at (0.5,0.25)  {$\times$};
  \node[class1, font=\tiny] at (0.9,0.25)  {$\times$};
  \node[class1, font=\tiny] at (0.3,0.45)  {$\times$};
  \node[class1, font=\tiny] at (1.1,0.45)  {$\times$};
  % ○ (class 0) -- above the parabola, spread across upper region
  \node[class0, font=\tiny] at (0.7,0.9)   {$\circ$};
  \node[class0, font=\tiny] at (0.35,0.85) {$\circ$};
  \node[class0, font=\tiny] at (1.05,0.85) {$\circ$};
  \node[class0, font=\tiny] at (0.1,1.3)   {$\circ$};
  \node[class0, font=\tiny] at (1.3,1.3)   {$\circ$};
  \node[class0, font=\tiny] at (0.7,1.2)   {$\circ$};
  \draw[testred, thick] plot[domain=0:1.4, samples=60] (\x, {2*(\x-0.7)*(\x-0.7)+0.35});
}}

% Wiggly separation: an irregular oscillating line cuts the box into two
% half-spaces; ○ fills the upper half, × fills the lower half.
% HOW TO ADJUST:
%   - The wiggly line is "plot[smooth] coordinates {...}": move waypoints to
%     change the frequency/amplitude of the oscillations.  The line oscillates
%     between y≈0.52 and y≈0.82, so keep × points below y≈0.4 and ○ points
%     above y≈0.95 to ensure they stay visually on the correct side.
%   - Add/remove nodes to adjust point density in each half.
\newcommand{\hypWiggly}{\hypIconBase{
  % half-space fills (clipped to box by \hypIconBase)
  \fill[class0!25] plot[smooth] coordinates {(-0.1,0.65) (0.2,0.82) (0.45,0.52) (0.7,0.78) (0.95,0.55) (1.2,0.72) (1.5,0.62)} -- (1.55,1.55) -- (-0.15,1.55) -- cycle;
  \fill[class1!25] plot[smooth] coordinates {(-0.1,0.65) (0.2,0.82) (0.45,0.52) (0.7,0.78) (0.95,0.55) (1.2,0.72) (1.5,0.62)} -- (1.55,-0.1) -- (-0.15,-0.1) -- cycle;
  % ○ (class 0) -- above wiggly line, spread across upper half-space
  \node[class0, font=\tiny] at (0.15,1.25) {$\circ$};
  \node[class0, font=\tiny] at (0.55,1.1)  {$\circ$};
  \node[class0, font=\tiny] at (1.05,1.2)  {$\circ$};
  \node[class0, font=\tiny] at (0.35,1.3)  {$\circ$};
  \node[class0, font=\tiny] at (0.85,1.1)  {$\circ$};
  \node[class0, font=\tiny] at (1.3,1.05)  {$\circ$};
  % × (class 1) -- below wiggly line, spread across lower half-space
  \node[class1, font=\tiny] at (0.15,0.22) {$\times$};
  \node[class1, font=\tiny] at (0.55,0.15) {$\times$};
  \node[class1, font=\tiny] at (1.05,0.25) {$\times$};
  \node[class1, font=\tiny] at (0.35,0.32) {$\times$};
  \node[class1, font=\tiny] at (0.85,0.2)  {$\times$};
  \node[class1, font=\tiny] at (1.3,0.35)  {$\times$};
  \draw[testred, thick] plot[smooth] coordinates {(-0.1,0.65) (0.2,0.82) (0.45,0.52) (0.7,0.78) (0.95,0.55) (1.2,0.72) (1.5,0.62)};
}}

% One "hypothesis space histogram" panel (used twice on the Meta Level
% slide: once for the prior $\Phi(\varphi)$, once for the posterior
% $\Phi(\varphi \mid D_n)$) -- a bar over $\varphi$ per representative
% hypothesis, with its decision-boundary icon underneath.
%
% Usage: \metapanel{<y-axis label>}{<bar1 height>}{<bar2>}{<bar3>}{<bar4>}
%   e.g. \metapanel{$\Phi(\varphi)$}{0.9}{1.0}{0.85}{0.8}
%
% HOW TO ADJUST:
%   - #1 = the y-axis label (put $\Phi(\varphi)$ or $\Phi(\varphi\mid
%     D_n)$ here when calling it).
%   - #2..#5 = the four bar heights, in that order; these are the ONLY
%     thing that differs between the "prior" and "posterior" calls on the
%     Meta Level slide -- raise #3 (hypLinB's bar), say, to show the
%     posterior concentrating on that hypothesis.
%   - Each bar/icon pair is one hardcoded block: "rectangle (x0,0)
%     rectangle (x1,height)" draws the bar, "\node at (xc,-0.85)
%     {\hyp...}" places its icon below the axis. To add a 5th hypothesis,
%     copy one block, shift its x-coordinates past 8.65, extend the axis
%     \draw's endpoint (currently 8.8), and add a 6th macro argument.
%   - "$\cdots$" nodes are just the ellipsis marks between bars,
%     indicating "more hypotheses in between" -- move them if you resize
%     the gaps.
%   - "scale=0.5" shrinks the whole panel to fit two of them side by side
%     on one slide; raise/lower it (together with the minipage height and
%     column widths in main.tex) if you change how many panels you show.
\newcommand{\metapanel}[5]{%
  \begin{tikzpicture}[scale=0.5]
    \draw[-{Latex}] (-1.9,0) -- (10,0) node[right] {$\varphi$};
    \draw[-{Latex}] (-1.9,-0.2) -- (-1.9,3.7) node[above] {#1};
    
    \draw[fill=testred!70, draw=testred] (-1.2,0) rectangle (0.2,#2);
    \node at (-0.5,-1.85) {\scalebox{0.8}{\hypLinA}};
    \node at (1.0,0.35) {$\cdots$};
    
    \draw[fill=testred!70, draw=testred] (1.8,0) rectangle (3.2,#3);
    \node at (2.5,-1.85) {\scalebox{0.8}{\hypLinB}};
    \node at (4.05,0.35) {$\cdots$};
    
    \draw[fill=testred!70, draw=testred] (4.8,0) rectangle (6.2,#4);
    \node at (5.5,-1.85) {\scalebox{0.8}{\hypCurve}};
    \node at (7.1,0.35) {$\cdots$};
    
    \draw[fill=testred!70, draw=testred] (7.8,0) rectangle (9.2,#5);
    \node at (8.5,-1.85) {\scalebox{0.8}{\hypWiggly}};
  \end{tikzpicture}%
}
.
% =================================================

% -----------------------------------------------
% Colors for feature-space / hypothesis-space sketches
% -----------------------------------------------
% class0/class1 = the two classes (o / x) everywhere below; hlA/hlB = the
% two highlight-ring colors used for x^(1)/x^(2) on the Statistical Model
% slide; testred = test points, decision-boundary curves, and histogram
% bars. Change these once here to re-theme every plot in the deck.
\definecolor{class0}{RGB}{56,142,60}    % green: y = 0, "o"
\definecolor{class1}{RGB}{81,45,168}    % purple: y = 1, "x"
\definecolor{hlA}{RGB}{216,27,96}       % pink highlight
\definecolor{hlB}{RGB}{230,81,0}        % orange highlight
\definecolor{testred}{RGB}{198,40,40}   % red: test point / boundary

% -----------------------------------------------
% Reusable TikZ pieces
% -----------------------------------------------
% Background point cloud shared by the two feature-space plots (Statistical
% Model / SL Framework frames): a class-0 ("o") cluster top-right, a
% class-1 ("x") cluster bottom-left.
%
% HOW TO ADJUST:
%   - The two \fill ellipses are just the soft shaded "cloud" background;
%     their (center) and (radius and radius) control where each cluster
%     sits and how big/spread-out it looks. They don't need to track the
%     individual points below -- they're a backdrop, not a fit.
%   - Each \node ... {$\circ$} / {$\times$} is ONE data point: the
%     coordinate after "at" is its (x2,x1) position (the plot has x2
%     horizontal, x1 vertical -- see the axis \draw's below). Add, remove,
%     or move these to change what D_n "looks like".
%   - The two \draw[-{Latex}] lines are the axes: change the second
%     coordinate to lengthen/shorten an axis; change the node[right]/
%     [above] text to relabel it.
%   - If you move anything outside roughly x2,x1 in [-0.5,6.4] x
%     [-0.5,4.6], also widen \fixedbbox below, or the picture will
%     silently resize and the plot will start "jumping" between slides
%     again (see \fixedbbox's own comment).
\newcommand{\featurecloud}{
  \fill[class0!20] (2.6,2.7) ellipse (1.35 and 1.05);
  \fill[class1!20] (1.0,1.0) ellipse (1.35 and 1.05);
  \node[class0] at (2.2,3.1) {$\circ$};
  \node[class0] at (2.6,3.35) {$\circ$};
  \node[class0] at (3.05,3.05) {$\circ$};
  \node[class0] at (2.35,2.65) {$\circ$};
  \node[class0] at (2.95,2.5) {$\circ$};
  \node[class0] at (3.35,2.85) {$\circ$};
  \node[class1] at (0.5,1.4) {$\times$};
  \node[class1] at (0.8,0.8) {$\times$};
  \node[class1] at (1.15,1.6) {$\times$};
  \node[class1] at (0.4,0.6) {$\times$};
  \node[class1] at (1.3,1.0) {$\times$};
  \node[class1] at (0.95,1.9) {$\times$};
  \draw[-{Latex}] (-0.3,0) -- (4.3,0) node[right] {$x_2$};
  \draw[-{Latex}] (0,-0.3) -- (0,4.3) node[above] {$x_1$};
}

% Small "$y=0$ / $y=1$" key box placed next to \featurecloud.
% HOW TO ADJUST:
%   - "5.1" in "(5.1+#1,1.7)" is the box's horizontal position (push it
%     further right if you widen \featurecloud, e.g. add more points out
%     past x2=4.3); "1.7" is its vertical position; #1 is an optional
%     extra offset a caller can pass, e.g. \featurelegend[0.3], without
%     editing this macro.
%   - "inner sep=4pt" / "font=\footnotesize" control the box's padding
%     and text size.
\newcommand{\featurelegend}[1][0]{
  \node[draw, align=left, font=\footnotesize, fill=white, inner sep=4pt]
    at (5.1+#1,1.7) {{\color{class0}$\circ$}\; $y=0$ \\ {\color{class1}$\times$}\; $y=1$};
}

% Invisible path that pins the bounding box of the feature-space plots to a
% fixed rectangle, so the picture is the same size on every slide that
% uses it (Statistical Model / SL Framework) regardless of which extra
% annotations (highlight circles, the test star, ...) are drawn on top --
% this is what keeps the plot from visually jumping when you flip between
% those slides.
%
% HOW TO ADJUST:
%   - Arguments are (xmin,ymin) rectangle (xmax,ymax), same (x2,x1)
%     coordinates as \featurecloud. Call this FIRST inside every
%     tikzpicture that uses \featurecloud, so all of them agree on one
%     size.
%   - If a new annotation you add (a star, a label, ...) sits outside
%     this rectangle, the picture's natural bounding box will grow to fit
%     it on THAT slide only, breaking the "same position everywhere"
%     guarantee. Widen the rectangle here (once, for every slide) instead
%     of nudging it per-frame.
\newcommand{\fixedbbox}{\path (-0.5,-0.5) rectangle (6.4,4.6);}

% "Dataset + decision boundary" icons used on the hypothesis-space slides
% (Meta Level bar charts, and the inline $\varphi_1,\varphi_2$ icons on
% the PPD slide): each hypothesis $\varphi$ is a *different* data
% generation mechanism, so each icon below has its own point-cloud
% geometry; the red curve just traces the shape of that particular cloud
% rather than being an independent line drawn over one fixed background.
%
% HOW TO ADJUST (applies to \hypIconBase and all four icons below it):
%   - All coordinates live in a fixed 1.4 x 1.4 box (see the \clip and
%     the border \draw here) -- keep new content roughly inside
%     [0,1.4] x [0,1.4] or it gets clipped off.
%   - #1 is the icon-specific content (fills, points, boundary curve),
%     supplied by each \hyp... macro below.
%   - These icons are always used heavily scaled down (via \scalebox or
%     inside \metapanel), so keep shapes simple/high-contrast -- fine
%     detail won't survive the shrink.
\newcommand{\hypIconBase}[1]{%
  \begin{tikzpicture}[baseline=(current bounding box.center)]
    \clip (0,0) rectangle (1.4,1.4);
    #1
    \draw (0,0) rectangle (1.4,1.4);
  \end{tikzpicture}%
}

% Main-diagonal split: "o" cluster top-right, "x" cluster bottom-left.
% HOW TO ADJUST:
%   - The two \fill ellipses are this hypothesis's cluster "shape"
%     (center + radii) -- move/resize them to change what the data looks
%     like under this $\varphi$.
%   - Each \node {$\circ$}/{$\times$} is one visible point; the final
%     \draw is the decision-boundary line -- its endpoints should track
%     wherever you put the two clusters (i.e. redraw it if you move the
%     ellipses, so it still separates them).
\newcommand{\hypLinA}{\hypIconBase{
  % half-space fills (clipped to box by \hypIconBase)
  \fill[class0!25] (-0.15,1.25) -- (1.55,0.15) -- (1.55,1.55) -- (-0.15,1.55) -- cycle;
  \fill[class1!25] (-0.15,1.25) -- (1.55,0.15) -- (1.55,-0.1) -- (-0.15,-0.1) -- cycle;
  \node[class0, font=\tiny] at (0.8,1.15) {$\circ$};
  \node[class0, font=\tiny] at (1.15,1.05) {$\circ$};
  \node[class0, font=\tiny] at (0.95,0.85) {$\circ$};
  \node[class1, font=\tiny] at (0.3,0.5) {$\times$};
  \node[class1, font=\tiny] at (0.6,0.3) {$\times$};
  \node[class1, font=\tiny] at (0.35,0.25) {$\times$};
  \draw[testred, thick] (-0.15,1.25) -- (1.55,0.15);
}}

% Anti-diagonal split: "x" cluster top-left, "o" cluster bottom-right --
% a genuinely different data layout from \hypLinA, not just a new line.
% HOW TO ADJUST: same recipe as \hypLinA (ellipses = clusters, nodes =
% points, final \draw = boundary line through the gap between them).
\newcommand{\hypLinB}{\hypIconBase{
  % half-space fills (clipped to box by \hypIconBase)
  \fill[class1!25] (-0.15,0.15) -- (1.55,1.25) -- (1.55,1.55) -- (-0.15,1.55) -- cycle;
  \fill[class0!25] (-0.15,0.15) -- (1.55,1.25) -- (1.55,-0.1) -- (-0.15,-0.1) -- cycle;
  \node[class1, font=\tiny] at (0.3,1.15) {$\times$};
  \node[class1, font=\tiny] at (0.6,1.05) {$\times$};
  \node[class1, font=\tiny] at (0.4,0.85) {$\times$};
  \node[class0, font=\tiny] at (1.1,0.55) {$\circ$};
  \node[class0, font=\tiny] at (0.85,0.25) {$\circ$};
  \node[class0, font=\tiny] at (1.15,0.2) {$\circ$};
  \draw[testred, thick] (-0.15,0.15) -- (1.55,1.25);
}}

% Quadratic split (y ~ x^2): × fills the bottom valley, ○ fills above.
% Boundary is the exact parabola y = 2(x-0.7)^2 + 0.35 (vertex at (0.7,0.35),
% arms reach y≈1.33 at the sides of the 1.4×1.4 box).
% HOW TO ADJUST:
%   - Vertex position: change (0.7) and (0.35) in the formula.
%   - Curvature: change the leading coefficient (currently 2).
%   - Data points: × must satisfy y < 2(x-0.7)^2+0.35; ○ must be above.
\newcommand{\hypCurve}{\hypIconBase{
  % half-space fills using the exact parabola (clipped to box by \hypIconBase)
  \fill[class0!25] plot[domain=-0.15:1.55, samples=60] (\x, {2*(\x-0.7)*(\x-0.7)+0.35})
    -- (1.55,1.55) -- (-0.15,1.55) -- cycle;
  \fill[class1!25] plot[domain=-0.15:1.55, samples=60] (\x, {2*(\x-0.7)*(\x-0.7)+0.35})
    -- (1.55,-0.1) -- (-0.15,-0.1) -- cycle;
  % × (class 1) -- below the parabola, spread across bottom valley
  \node[class1, font=\tiny] at (0.7,0.15)  {$\times$};
  \node[class1, font=\tiny] at (0.5,0.25)  {$\times$};
  \node[class1, font=\tiny] at (0.9,0.25)  {$\times$};
  \node[class1, font=\tiny] at (0.3,0.45)  {$\times$};
  \node[class1, font=\tiny] at (1.1,0.45)  {$\times$};
  % ○ (class 0) -- above the parabola, spread across upper region
  \node[class0, font=\tiny] at (0.7,0.9)   {$\circ$};
  \node[class0, font=\tiny] at (0.35,0.85) {$\circ$};
  \node[class0, font=\tiny] at (1.05,0.85) {$\circ$};
  \node[class0, font=\tiny] at (0.1,1.3)   {$\circ$};
  \node[class0, font=\tiny] at (1.3,1.3)   {$\circ$};
  \node[class0, font=\tiny] at (0.7,1.2)   {$\circ$};
  \draw[testred, thick] plot[domain=0:1.4, samples=60] (\x, {2*(\x-0.7)*(\x-0.7)+0.35});
}}

% Wiggly separation: an irregular oscillating line cuts the box into two
% half-spaces; ○ fills the upper half, × fills the lower half.
% HOW TO ADJUST:
%   - The wiggly line is "plot[smooth] coordinates {...}": move waypoints to
%     change the frequency/amplitude of the oscillations.  The line oscillates
%     between y≈0.52 and y≈0.82, so keep × points below y≈0.4 and ○ points
%     above y≈0.95 to ensure they stay visually on the correct side.
%   - Add/remove nodes to adjust point density in each half.
\newcommand{\hypWiggly}{\hypIconBase{
  % half-space fills (clipped to box by \hypIconBase)
  \fill[class0!25] plot[smooth] coordinates {(-0.1,0.65) (0.2,0.82) (0.45,0.52) (0.7,0.78) (0.95,0.55) (1.2,0.72) (1.5,0.62)} -- (1.55,1.55) -- (-0.15,1.55) -- cycle;
  \fill[class1!25] plot[smooth] coordinates {(-0.1,0.65) (0.2,0.82) (0.45,0.52) (0.7,0.78) (0.95,0.55) (1.2,0.72) (1.5,0.62)} -- (1.55,-0.1) -- (-0.15,-0.1) -- cycle;
  % ○ (class 0) -- above wiggly line, spread across upper half-space
  \node[class0, font=\tiny] at (0.15,1.25) {$\circ$};
  \node[class0, font=\tiny] at (0.55,1.1)  {$\circ$};
  \node[class0, font=\tiny] at (1.05,1.2)  {$\circ$};
  \node[class0, font=\tiny] at (0.35,1.3)  {$\circ$};
  \node[class0, font=\tiny] at (0.85,1.1)  {$\circ$};
  \node[class0, font=\tiny] at (1.3,1.05)  {$\circ$};
  % × (class 1) -- below wiggly line, spread across lower half-space
  \node[class1, font=\tiny] at (0.15,0.22) {$\times$};
  \node[class1, font=\tiny] at (0.55,0.15) {$\times$};
  \node[class1, font=\tiny] at (1.05,0.25) {$\times$};
  \node[class1, font=\tiny] at (0.35,0.32) {$\times$};
  \node[class1, font=\tiny] at (0.85,0.2)  {$\times$};
  \node[class1, font=\tiny] at (1.3,0.35)  {$\times$};
  \draw[testred, thick] plot[smooth] coordinates {(-0.1,0.65) (0.2,0.82) (0.45,0.52) (0.7,0.78) (0.95,0.55) (1.2,0.72) (1.5,0.62)};
}}

% One "hypothesis space histogram" panel (used twice on the Meta Level
% slide: once for the prior $\Phi(\varphi)$, once for the posterior
% $\Phi(\varphi \mid D_n)$) -- a bar over $\varphi$ per representative
% hypothesis, with its decision-boundary icon underneath.
%
% Usage: \metapanel{<y-axis label>}{<bar1 height>}{<bar2>}{<bar3>}{<bar4>}
%   e.g. \metapanel{$\Phi(\varphi)$}{0.9}{1.0}{0.85}{0.8}
%
% HOW TO ADJUST:
%   - #1 = the y-axis label (put $\Phi(\varphi)$ or $\Phi(\varphi\mid
%     D_n)$ here when calling it).
%   - #2..#5 = the four bar heights, in that order; these are the ONLY
%     thing that differs between the "prior" and "posterior" calls on the
%     Meta Level slide -- raise #3 (hypLinB's bar), say, to show the
%     posterior concentrating on that hypothesis.
%   - Each bar/icon pair is one hardcoded block: "rectangle (x0,0)
%     rectangle (x1,height)" draws the bar, "\node at (xc,-0.85)
%     {\hyp...}" places its icon below the axis. To add a 5th hypothesis,
%     copy one block, shift its x-coordinates past 8.65, extend the axis
%     \draw's endpoint (currently 8.8), and add a 6th macro argument.
%   - "$\cdots$" nodes are just the ellipsis marks between bars,
%     indicating "more hypotheses in between" -- move them if you resize
%     the gaps.
%   - "scale=0.5" shrinks the whole panel to fit two of them side by side
%     on one slide; raise/lower it (together with the minipage height and
%     column widths in main.tex) if you change how many panels you show.
\newcommand{\metapanel}[5]{%
  \begin{tikzpicture}[scale=0.5]
    \draw[-{Latex}] (-1.9,0) -- (10,0) node[right] {$\varphi$};
    \draw[-{Latex}] (-1.9,-0.2) -- (-1.9,3.7) node[above] {#1};
    
    \draw[fill=testred!70, draw=testred] (-1.2,0) rectangle (0.2,#2);
    \node at (-0.5,-1.85) {\scalebox{0.8}{\hypLinA}};
    \node at (1.0,0.35) {$\cdots$};
    
    \draw[fill=testred!70, draw=testred] (1.8,0) rectangle (3.2,#3);
    \node at (2.5,-1.85) {\scalebox{0.8}{\hypLinB}};
    \node at (4.05,0.35) {$\cdots$};
    
    \draw[fill=testred!70, draw=testred] (4.8,0) rectangle (6.2,#4);
    \node at (5.5,-1.85) {\scalebox{0.8}{\hypCurve}};
    \node at (7.1,0.35) {$\cdots$};
    
    \draw[fill=testred!70, draw=testred] (7.8,0) rectangle (9.2,#5);
    \node at (8.5,-1.85) {\scalebox{0.8}{\hypWiggly}};
  \end{tikzpicture}%
}
.
% =================================================

% -----------------------------------------------
% Colors for feature-space / hypothesis-space sketches
% -----------------------------------------------
% class0/class1 = the two classes (o / x) everywhere below; hlA/hlB = the
% two highlight-ring colors used for x^(1)/x^(2) on the Statistical Model
% slide; testred = test points, decision-boundary curves, and histogram
% bars. Change these once here to re-theme every plot in the deck.
\definecolor{class0}{RGB}{56,142,60}    % green: y = 0, "o"
\definecolor{class1}{RGB}{81,45,168}    % purple: y = 1, "x"
\definecolor{hlA}{RGB}{216,27,96}       % pink highlight
\definecolor{hlB}{RGB}{230,81,0}        % orange highlight
\definecolor{testred}{RGB}{198,40,40}   % red: test point / boundary

% -----------------------------------------------
% Reusable TikZ pieces
% -----------------------------------------------
% Background point cloud shared by the two feature-space plots (Statistical
% Model / SL Framework frames): a class-0 ("o") cluster top-right, a
% class-1 ("x") cluster bottom-left.
%
% HOW TO ADJUST:
%   - The two \fill ellipses are just the soft shaded "cloud" background;
%     their (center) and (radius and radius) control where each cluster
%     sits and how big/spread-out it looks. They don't need to track the
%     individual points below -- they're a backdrop, not a fit.
%   - Each \node ... {$\circ$} / {$\times$} is ONE data point: the
%     coordinate after "at" is its (x2,x1) position (the plot has x2
%     horizontal, x1 vertical -- see the axis \draw's below). Add, remove,
%     or move these to change what D_n "looks like".
%   - The two \draw[-{Latex}] lines are the axes: change the second
%     coordinate to lengthen/shorten an axis; change the node[right]/
%     [above] text to relabel it.
%   - If you move anything outside roughly x2,x1 in [-0.5,6.4] x
%     [-0.5,4.6], also widen \fixedbbox below, or the picture will
%     silently resize and the plot will start "jumping" between slides
%     again (see \fixedbbox's own comment).
\newcommand{\featurecloud}{
  \fill[class0!20] (2.6,2.7) ellipse (1.35 and 1.05);
  \fill[class1!20] (1.0,1.0) ellipse (1.35 and 1.05);
  \node[class0] at (2.2,3.1) {$\circ$};
  \node[class0] at (2.6,3.35) {$\circ$};
  \node[class0] at (3.05,3.05) {$\circ$};
  \node[class0] at (2.35,2.65) {$\circ$};
  \node[class0] at (2.95,2.5) {$\circ$};
  \node[class0] at (3.35,2.85) {$\circ$};
  \node[class1] at (0.5,1.4) {$\times$};
  \node[class1] at (0.8,0.8) {$\times$};
  \node[class1] at (1.15,1.6) {$\times$};
  \node[class1] at (0.4,0.6) {$\times$};
  \node[class1] at (1.3,1.0) {$\times$};
  \node[class1] at (0.95,1.9) {$\times$};
  \draw[-{Latex}] (-0.3,0) -- (4.3,0) node[right] {$x_2$};
  \draw[-{Latex}] (0,-0.3) -- (0,4.3) node[above] {$x_1$};
}

% Small "$y=0$ / $y=1$" key box placed next to \featurecloud.
% HOW TO ADJUST:
%   - "5.1" in "(5.1+#1,1.7)" is the box's horizontal position (push it
%     further right if you widen \featurecloud, e.g. add more points out
%     past x2=4.3); "1.7" is its vertical position; #1 is an optional
%     extra offset a caller can pass, e.g. \featurelegend[0.3], without
%     editing this macro.
%   - "inner sep=4pt" / "font=\footnotesize" control the box's padding
%     and text size.
\newcommand{\featurelegend}[1][0]{
  \node[draw, align=left, font=\footnotesize, fill=white, inner sep=4pt]
    at (5.1+#1,1.7) {{\color{class0}$\circ$}\; $y=0$ \\ {\color{class1}$\times$}\; $y=1$};
}

% Invisible path that pins the bounding box of the feature-space plots to a
% fixed rectangle, so the picture is the same size on every slide that
% uses it (Statistical Model / SL Framework) regardless of which extra
% annotations (highlight circles, the test star, ...) are drawn on top --
% this is what keeps the plot from visually jumping when you flip between
% those slides.
%
% HOW TO ADJUST:
%   - Arguments are (xmin,ymin) rectangle (xmax,ymax), same (x2,x1)
%     coordinates as \featurecloud. Call this FIRST inside every
%     tikzpicture that uses \featurecloud, so all of them agree on one
%     size.
%   - If a new annotation you add (a star, a label, ...) sits outside
%     this rectangle, the picture's natural bounding box will grow to fit
%     it on THAT slide only, breaking the "same position everywhere"
%     guarantee. Widen the rectangle here (once, for every slide) instead
%     of nudging it per-frame.
\newcommand{\fixedbbox}{\path (-0.5,-0.5) rectangle (6.4,4.6);}

% "Dataset + decision boundary" icons used on the hypothesis-space slides
% (Meta Level bar charts, and the inline $\varphi_1,\varphi_2$ icons on
% the PPD slide): each hypothesis $\varphi$ is a *different* data
% generation mechanism, so each icon below has its own point-cloud
% geometry; the red curve just traces the shape of that particular cloud
% rather than being an independent line drawn over one fixed background.
%
% HOW TO ADJUST (applies to \hypIconBase and all four icons below it):
%   - All coordinates live in a fixed 1.4 x 1.4 box (see the \clip and
%     the border \draw here) -- keep new content roughly inside
%     [0,1.4] x [0,1.4] or it gets clipped off.
%   - #1 is the icon-specific content (fills, points, boundary curve),
%     supplied by each \hyp... macro below.
%   - These icons are always used heavily scaled down (via \scalebox or
%     inside \metapanel), so keep shapes simple/high-contrast -- fine
%     detail won't survive the shrink.
\newcommand{\hypIconBase}[1]{%
  \begin{tikzpicture}[baseline=(current bounding box.center)]
    \clip (0,0) rectangle (1.4,1.4);
    #1
    \draw (0,0) rectangle (1.4,1.4);
  \end{tikzpicture}%
}

% Main-diagonal split: "o" cluster top-right, "x" cluster bottom-left.
% HOW TO ADJUST:
%   - The two \fill ellipses are this hypothesis's cluster "shape"
%     (center + radii) -- move/resize them to change what the data looks
%     like under this $\varphi$.
%   - Each \node {$\circ$}/{$\times$} is one visible point; the final
%     \draw is the decision-boundary line -- its endpoints should track
%     wherever you put the two clusters (i.e. redraw it if you move the
%     ellipses, so it still separates them).
\newcommand{\hypLinA}{\hypIconBase{
  % half-space fills (clipped to box by \hypIconBase)
  \fill[class0!25] (-0.15,1.25) -- (1.55,0.15) -- (1.55,1.55) -- (-0.15,1.55) -- cycle;
  \fill[class1!25] (-0.15,1.25) -- (1.55,0.15) -- (1.55,-0.1) -- (-0.15,-0.1) -- cycle;
  \node[class0, font=\tiny] at (0.8,1.15) {$\circ$};
  \node[class0, font=\tiny] at (1.15,1.05) {$\circ$};
  \node[class0, font=\tiny] at (0.95,0.85) {$\circ$};
  \node[class1, font=\tiny] at (0.3,0.5) {$\times$};
  \node[class1, font=\tiny] at (0.6,0.3) {$\times$};
  \node[class1, font=\tiny] at (0.35,0.25) {$\times$};
  \draw[testred, thick] (-0.15,1.25) -- (1.55,0.15);
}}

% Anti-diagonal split: "x" cluster top-left, "o" cluster bottom-right --
% a genuinely different data layout from \hypLinA, not just a new line.
% HOW TO ADJUST: same recipe as \hypLinA (ellipses = clusters, nodes =
% points, final \draw = boundary line through the gap between them).
\newcommand{\hypLinB}{\hypIconBase{
  % half-space fills (clipped to box by \hypIconBase)
  \fill[class1!25] (-0.15,0.15) -- (1.55,1.25) -- (1.55,1.55) -- (-0.15,1.55) -- cycle;
  \fill[class0!25] (-0.15,0.15) -- (1.55,1.25) -- (1.55,-0.1) -- (-0.15,-0.1) -- cycle;
  \node[class1, font=\tiny] at (0.3,1.15) {$\times$};
  \node[class1, font=\tiny] at (0.6,1.05) {$\times$};
  \node[class1, font=\tiny] at (0.4,0.85) {$\times$};
  \node[class0, font=\tiny] at (1.1,0.55) {$\circ$};
  \node[class0, font=\tiny] at (0.85,0.25) {$\circ$};
  \node[class0, font=\tiny] at (1.15,0.2) {$\circ$};
  \draw[testred, thick] (-0.15,0.15) -- (1.55,1.25);
}}

% Quadratic split (y ~ x^2): × fills the bottom valley, ○ fills above.
% Boundary is the exact parabola y = 2(x-0.7)^2 + 0.35 (vertex at (0.7,0.35),
% arms reach y≈1.33 at the sides of the 1.4×1.4 box).
% HOW TO ADJUST:
%   - Vertex position: change (0.7) and (0.35) in the formula.
%   - Curvature: change the leading coefficient (currently 2).
%   - Data points: × must satisfy y < 2(x-0.7)^2+0.35; ○ must be above.
\newcommand{\hypCurve}{\hypIconBase{
  % half-space fills using the exact parabola (clipped to box by \hypIconBase)
  \fill[class0!25] plot[domain=-0.15:1.55, samples=60] (\x, {2*(\x-0.7)*(\x-0.7)+0.35})
    -- (1.55,1.55) -- (-0.15,1.55) -- cycle;
  \fill[class1!25] plot[domain=-0.15:1.55, samples=60] (\x, {2*(\x-0.7)*(\x-0.7)+0.35})
    -- (1.55,-0.1) -- (-0.15,-0.1) -- cycle;
  % × (class 1) -- below the parabola, spread across bottom valley
  \node[class1, font=\tiny] at (0.7,0.15)  {$\times$};
  \node[class1, font=\tiny] at (0.5,0.25)  {$\times$};
  \node[class1, font=\tiny] at (0.9,0.25)  {$\times$};
  \node[class1, font=\tiny] at (0.3,0.45)  {$\times$};
  \node[class1, font=\tiny] at (1.1,0.45)  {$\times$};
  % ○ (class 0) -- above the parabola, spread across upper region
  \node[class0, font=\tiny] at (0.7,0.9)   {$\circ$};
  \node[class0, font=\tiny] at (0.35,0.85) {$\circ$};
  \node[class0, font=\tiny] at (1.05,0.85) {$\circ$};
  \node[class0, font=\tiny] at (0.1,1.3)   {$\circ$};
  \node[class0, font=\tiny] at (1.3,1.3)   {$\circ$};
  \node[class0, font=\tiny] at (0.7,1.2)   {$\circ$};
  \draw[testred, thick] plot[domain=0:1.4, samples=60] (\x, {2*(\x-0.7)*(\x-0.7)+0.35});
}}

% Wiggly separation: an irregular oscillating line cuts the box into two
% half-spaces; ○ fills the upper half, × fills the lower half.
% HOW TO ADJUST:
%   - The wiggly line is "plot[smooth] coordinates {...}": move waypoints to
%     change the frequency/amplitude of the oscillations.  The line oscillates
%     between y≈0.52 and y≈0.82, so keep × points below y≈0.4 and ○ points
%     above y≈0.95 to ensure they stay visually on the correct side.
%   - Add/remove nodes to adjust point density in each half.
\newcommand{\hypWiggly}{\hypIconBase{
  % half-space fills (clipped to box by \hypIconBase)
  \fill[class0!25] plot[smooth] coordinates {(-0.1,0.65) (0.2,0.82) (0.45,0.52) (0.7,0.78) (0.95,0.55) (1.2,0.72) (1.5,0.62)} -- (1.55,1.55) -- (-0.15,1.55) -- cycle;
  \fill[class1!25] plot[smooth] coordinates {(-0.1,0.65) (0.2,0.82) (0.45,0.52) (0.7,0.78) (0.95,0.55) (1.2,0.72) (1.5,0.62)} -- (1.55,-0.1) -- (-0.15,-0.1) -- cycle;
  % ○ (class 0) -- above wiggly line, spread across upper half-space
  \node[class0, font=\tiny] at (0.15,1.25) {$\circ$};
  \node[class0, font=\tiny] at (0.55,1.1)  {$\circ$};
  \node[class0, font=\tiny] at (1.05,1.2)  {$\circ$};
  \node[class0, font=\tiny] at (0.35,1.3)  {$\circ$};
  \node[class0, font=\tiny] at (0.85,1.1)  {$\circ$};
  \node[class0, font=\tiny] at (1.3,1.05)  {$\circ$};
  % × (class 1) -- below wiggly line, spread across lower half-space
  \node[class1, font=\tiny] at (0.15,0.22) {$\times$};
  \node[class1, font=\tiny] at (0.55,0.15) {$\times$};
  \node[class1, font=\tiny] at (1.05,0.25) {$\times$};
  \node[class1, font=\tiny] at (0.35,0.32) {$\times$};
  \node[class1, font=\tiny] at (0.85,0.2)  {$\times$};
  \node[class1, font=\tiny] at (1.3,0.35)  {$\times$};
  \draw[testred, thick] plot[smooth] coordinates {(-0.1,0.65) (0.2,0.82) (0.45,0.52) (0.7,0.78) (0.95,0.55) (1.2,0.72) (1.5,0.62)};
}}

% One "hypothesis space histogram" panel (used twice on the Meta Level
% slide: once for the prior $\Phi(\varphi)$, once for the posterior
% $\Phi(\varphi \mid D_n)$) -- a bar over $\varphi$ per representative
% hypothesis, with its decision-boundary icon underneath.
%
% Usage: \metapanel{<y-axis label>}{<bar1 height>}{<bar2>}{<bar3>}{<bar4>}
%   e.g. \metapanel{$\Phi(\varphi)$}{0.9}{1.0}{0.85}{0.8}
%
% HOW TO ADJUST:
%   - #1 = the y-axis label (put $\Phi(\varphi)$ or $\Phi(\varphi\mid
%     D_n)$ here when calling it).
%   - #2..#5 = the four bar heights, in that order; these are the ONLY
%     thing that differs between the "prior" and "posterior" calls on the
%     Meta Level slide -- raise #3 (hypLinB's bar), say, to show the
%     posterior concentrating on that hypothesis.
%   - Each bar/icon pair is one hardcoded block: "rectangle (x0,0)
%     rectangle (x1,height)" draws the bar, "\node at (xc,-0.85)
%     {\hyp...}" places its icon below the axis. To add a 5th hypothesis,
%     copy one block, shift its x-coordinates past 8.65, extend the axis
%     \draw's endpoint (currently 8.8), and add a 6th macro argument.
%   - "$\cdots$" nodes are just the ellipsis marks between bars,
%     indicating "more hypotheses in between" -- move them if you resize
%     the gaps.
%   - "scale=0.5" shrinks the whole panel to fit two of them side by side
%     on one slide; raise/lower it (together with the minipage height and
%     column widths in main.tex) if you change how many panels you show.
\newcommand{\metapanel}[5]{%
  \begin{tikzpicture}[scale=0.5]
    \draw[-{Latex}] (-1.9,0) -- (10,0) node[right] {$\varphi$};
    \draw[-{Latex}] (-1.9,-0.2) -- (-1.9,3.7) node[above] {#1};
    
    \draw[fill=testred!70, draw=testred] (-1.2,0) rectangle (0.2,#2);
    \node at (-0.5,-1.85) {\scalebox{0.8}{\hypLinA}};
    \node at (1.0,0.35) {$\cdots$};
    
    \draw[fill=testred!70, draw=testred] (1.8,0) rectangle (3.2,#3);
    \node at (2.5,-1.85) {\scalebox{0.8}{\hypLinB}};
    \node at (4.05,0.35) {$\cdots$};
    
    \draw[fill=testred!70, draw=testred] (4.8,0) rectangle (6.2,#4);
    \node at (5.5,-1.85) {\scalebox{0.8}{\hypCurve}};
    \node at (7.1,0.35) {$\cdots$};
    
    \draw[fill=testred!70, draw=testred] (7.8,0) rectangle (9.2,#5);
    \node at (8.5,-1.85) {\scalebox{0.8}{\hypWiggly}};
  \end{tikzpicture}%
}
.
% =================================================

% -----------------------------------------------
% Colors for feature-space / hypothesis-space sketches
% -----------------------------------------------
% class0/class1 = the two classes (o / x) everywhere below; hlA/hlB = the
% two highlight-ring colors used for x^(1)/x^(2) on the Statistical Model
% slide; testred = test points, decision-boundary curves, and histogram
% bars. Change these once here to re-theme every plot in the deck.
\definecolor{class0}{RGB}{56,142,60}    % green: y = 0, "o"
\definecolor{class1}{RGB}{81,45,168}    % purple: y = 1, "x"
\definecolor{hlA}{RGB}{216,27,96}       % pink highlight
\definecolor{hlB}{RGB}{230,81,0}        % orange highlight
\definecolor{testred}{RGB}{198,40,40}   % red: test point / boundary

% -----------------------------------------------
% Reusable TikZ pieces
% -----------------------------------------------
% Background point cloud shared by the two feature-space plots (Statistical
% Model / SL Framework frames): a class-0 ("o") cluster top-right, a
% class-1 ("x") cluster bottom-left.
%
% HOW TO ADJUST:
%   - The two \fill ellipses are just the soft shaded "cloud" background;
%     their (center) and (radius and radius) control where each cluster
%     sits and how big/spread-out it looks. They don't need to track the
%     individual points below -- they're a backdrop, not a fit.
%   - Each \node ... {$\circ$} / {$\times$} is ONE data point: the
%     coordinate after "at" is its (x2,x1) position (the plot has x2
%     horizontal, x1 vertical -- see the axis \draw's below). Add, remove,
%     or move these to change what D_n "looks like".
%   - The two \draw[-{Latex}] lines are the axes: change the second
%     coordinate to lengthen/shorten an axis; change the node[right]/
%     [above] text to relabel it.
%   - If you move anything outside roughly x2,x1 in [-0.5,6.4] x
%     [-0.5,4.6], also widen \fixedbbox below, or the picture will
%     silently resize and the plot will start "jumping" between slides
%     again (see \fixedbbox's own comment).
\newcommand{\featurecloud}{
  \fill[class0!20] (2.6,2.7) ellipse (1.35 and 1.05);
  \fill[class1!20] (1.0,1.0) ellipse (1.35 and 1.05);
  \node[class0] at (2.2,3.1) {$\circ$};
  \node[class0] at (2.6,3.35) {$\circ$};
  \node[class0] at (3.05,3.05) {$\circ$};
  \node[class0] at (2.35,2.65) {$\circ$};
  \node[class0] at (2.95,2.5) {$\circ$};
  \node[class0] at (3.35,2.85) {$\circ$};
  \node[class1] at (0.5,1.4) {$\times$};
  \node[class1] at (0.8,0.8) {$\times$};
  \node[class1] at (1.15,1.6) {$\times$};
  \node[class1] at (0.4,0.6) {$\times$};
  \node[class1] at (1.3,1.0) {$\times$};
  \node[class1] at (0.95,1.9) {$\times$};
  \draw[-{Latex}] (-0.3,0) -- (4.3,0) node[right] {$x_2$};
  \draw[-{Latex}] (0,-0.3) -- (0,4.3) node[above] {$x_1$};
}

% Small "$y=0$ / $y=1$" key box placed next to \featurecloud.
% HOW TO ADJUST:
%   - "5.1" in "(5.1+#1,1.7)" is the box's horizontal position (push it
%     further right if you widen \featurecloud, e.g. add more points out
%     past x2=4.3); "1.7" is its vertical position; #1 is an optional
%     extra offset a caller can pass, e.g. \featurelegend[0.3], without
%     editing this macro.
%   - "inner sep=4pt" / "font=\footnotesize" control the box's padding
%     and text size.
\newcommand{\featurelegend}[1][0]{
  \node[draw, align=left, font=\footnotesize, fill=white, inner sep=4pt]
    at (5.1+#1,1.7) {{\color{class0}$\circ$}\; $y=0$ \\ {\color{class1}$\times$}\; $y=1$};
}

% Invisible path that pins the bounding box of the feature-space plots to a
% fixed rectangle, so the picture is the same size on every slide that
% uses it (Statistical Model / SL Framework) regardless of which extra
% annotations (highlight circles, the test star, ...) are drawn on top --
% this is what keeps the plot from visually jumping when you flip between
% those slides.
%
% HOW TO ADJUST:
%   - Arguments are (xmin,ymin) rectangle (xmax,ymax), same (x2,x1)
%     coordinates as \featurecloud. Call this FIRST inside every
%     tikzpicture that uses \featurecloud, so all of them agree on one
%     size.
%   - If a new annotation you add (a star, a label, ...) sits outside
%     this rectangle, the picture's natural bounding box will grow to fit
%     it on THAT slide only, breaking the "same position everywhere"
%     guarantee. Widen the rectangle here (once, for every slide) instead
%     of nudging it per-frame.
\newcommand{\fixedbbox}{\path (-0.5,-0.5) rectangle (6.4,4.6);}

% "Dataset + decision boundary" icons used on the hypothesis-space slides
% (Meta Level bar charts, and the inline $\varphi_1,\varphi_2$ icons on
% the PPD slide): each hypothesis $\varphi$ is a *different* data
% generation mechanism, so each icon below has its own point-cloud
% geometry; the red curve just traces the shape of that particular cloud
% rather than being an independent line drawn over one fixed background.
%
% HOW TO ADJUST (applies to \hypIconBase and all four icons below it):
%   - All coordinates live in a fixed 1.4 x 1.4 box (see the \clip and
%     the border \draw here) -- keep new content roughly inside
%     [0,1.4] x [0,1.4] or it gets clipped off.
%   - #1 is the icon-specific content (fills, points, boundary curve),
%     supplied by each \hyp... macro below.
%   - These icons are always used heavily scaled down (via \scalebox or
%     inside \metapanel), so keep shapes simple/high-contrast -- fine
%     detail won't survive the shrink.
\newcommand{\hypIconBase}[1]{%
  \begin{tikzpicture}[baseline=(current bounding box.center)]
    \clip (0,0) rectangle (1.4,1.4);
    #1
    \draw (0,0) rectangle (1.4,1.4);
  \end{tikzpicture}%
}

% Main-diagonal split: "o" cluster top-right, "x" cluster bottom-left.
% HOW TO ADJUST:
%   - The two \fill ellipses are this hypothesis's cluster "shape"
%     (center + radii) -- move/resize them to change what the data looks
%     like under this $\varphi$.
%   - Each \node {$\circ$}/{$\times$} is one visible point; the final
%     \draw is the decision-boundary line -- its endpoints should track
%     wherever you put the two clusters (i.e. redraw it if you move the
%     ellipses, so it still separates them).
\newcommand{\hypLinA}{\hypIconBase{
  % half-space fills (clipped to box by \hypIconBase)
  \fill[class0!25] (-0.15,1.25) -- (1.55,0.15) -- (1.55,1.55) -- (-0.15,1.55) -- cycle;
  \fill[class1!25] (-0.15,1.25) -- (1.55,0.15) -- (1.55,-0.1) -- (-0.15,-0.1) -- cycle;
  \node[class0, font=\tiny] at (0.8,1.15) {$\circ$};
  \node[class0, font=\tiny] at (1.15,1.05) {$\circ$};
  \node[class0, font=\tiny] at (0.95,0.85) {$\circ$};
  \node[class1, font=\tiny] at (0.3,0.5) {$\times$};
  \node[class1, font=\tiny] at (0.6,0.3) {$\times$};
  \node[class1, font=\tiny] at (0.35,0.25) {$\times$};
  \draw[testred, thick] (-0.15,1.25) -- (1.55,0.15);
}}

% Anti-diagonal split: "x" cluster top-left, "o" cluster bottom-right --
% a genuinely different data layout from \hypLinA, not just a new line.
% HOW TO ADJUST: same recipe as \hypLinA (ellipses = clusters, nodes =
% points, final \draw = boundary line through the gap between them).
\newcommand{\hypLinB}{\hypIconBase{
  % half-space fills (clipped to box by \hypIconBase)
  \fill[class1!25] (-0.15,0.15) -- (1.55,1.25) -- (1.55,1.55) -- (-0.15,1.55) -- cycle;
  \fill[class0!25] (-0.15,0.15) -- (1.55,1.25) -- (1.55,-0.1) -- (-0.15,-0.1) -- cycle;
  \node[class1, font=\tiny] at (0.3,1.15) {$\times$};
  \node[class1, font=\tiny] at (0.6,1.05) {$\times$};
  \node[class1, font=\tiny] at (0.4,0.85) {$\times$};
  \node[class0, font=\tiny] at (1.1,0.55) {$\circ$};
  \node[class0, font=\tiny] at (0.85,0.25) {$\circ$};
  \node[class0, font=\tiny] at (1.15,0.2) {$\circ$};
  \draw[testred, thick] (-0.15,0.15) -- (1.55,1.25);
}}

% Quadratic split (y ~ x^2): × fills the bottom valley, ○ fills above.
% Boundary is the exact parabola y = 2(x-0.7)^2 + 0.35 (vertex at (0.7,0.35),
% arms reach y≈1.33 at the sides of the 1.4×1.4 box).
% HOW TO ADJUST:
%   - Vertex position: change (0.7) and (0.35) in the formula.
%   - Curvature: change the leading coefficient (currently 2).
%   - Data points: × must satisfy y < 2(x-0.7)^2+0.35; ○ must be above.
\newcommand{\hypCurve}{\hypIconBase{
  % half-space fills using the exact parabola (clipped to box by \hypIconBase)
  \fill[class0!25] plot[domain=-0.15:1.55, samples=60] (\x, {2*(\x-0.7)*(\x-0.7)+0.35})
    -- (1.55,1.55) -- (-0.15,1.55) -- cycle;
  \fill[class1!25] plot[domain=-0.15:1.55, samples=60] (\x, {2*(\x-0.7)*(\x-0.7)+0.35})
    -- (1.55,-0.1) -- (-0.15,-0.1) -- cycle;
  % × (class 1) -- below the parabola, spread across bottom valley
  \node[class1, font=\tiny] at (0.7,0.15)  {$\times$};
  \node[class1, font=\tiny] at (0.5,0.25)  {$\times$};
  \node[class1, font=\tiny] at (0.9,0.25)  {$\times$};
  \node[class1, font=\tiny] at (0.3,0.45)  {$\times$};
  \node[class1, font=\tiny] at (1.1,0.45)  {$\times$};
  % ○ (class 0) -- above the parabola, spread across upper region
  \node[class0, font=\tiny] at (0.7,0.9)   {$\circ$};
  \node[class0, font=\tiny] at (0.35,0.85) {$\circ$};
  \node[class0, font=\tiny] at (1.05,0.85) {$\circ$};
  \node[class0, font=\tiny] at (0.1,1.3)   {$\circ$};
  \node[class0, font=\tiny] at (1.3,1.3)   {$\circ$};
  \node[class0, font=\tiny] at (0.7,1.2)   {$\circ$};
  \draw[testred, thick] plot[domain=0:1.4, samples=60] (\x, {2*(\x-0.7)*(\x-0.7)+0.35});
}}

% Wiggly separation: an irregular oscillating line cuts the box into two
% half-spaces; ○ fills the upper half, × fills the lower half.
% HOW TO ADJUST:
%   - The wiggly line is "plot[smooth] coordinates {...}": move waypoints to
%     change the frequency/amplitude of the oscillations.  The line oscillates
%     between y≈0.52 and y≈0.82, so keep × points below y≈0.4 and ○ points
%     above y≈0.95 to ensure they stay visually on the correct side.
%   - Add/remove nodes to adjust point density in each half.
\newcommand{\hypWiggly}{\hypIconBase{
  % half-space fills (clipped to box by \hypIconBase)
  \fill[class0!25] plot[smooth] coordinates {(-0.1,0.65) (0.2,0.82) (0.45,0.52) (0.7,0.78) (0.95,0.55) (1.2,0.72) (1.5,0.62)} -- (1.55,1.55) -- (-0.15,1.55) -- cycle;
  \fill[class1!25] plot[smooth] coordinates {(-0.1,0.65) (0.2,0.82) (0.45,0.52) (0.7,0.78) (0.95,0.55) (1.2,0.72) (1.5,0.62)} -- (1.55,-0.1) -- (-0.15,-0.1) -- cycle;
  % ○ (class 0) -- above wiggly line, spread across upper half-space
  \node[class0, font=\tiny] at (0.15,1.25) {$\circ$};
  \node[class0, font=\tiny] at (0.55,1.1)  {$\circ$};
  \node[class0, font=\tiny] at (1.05,1.2)  {$\circ$};
  \node[class0, font=\tiny] at (0.35,1.3)  {$\circ$};
  \node[class0, font=\tiny] at (0.85,1.1)  {$\circ$};
  \node[class0, font=\tiny] at (1.3,1.05)  {$\circ$};
  % × (class 1) -- below wiggly line, spread across lower half-space
  \node[class1, font=\tiny] at (0.15,0.22) {$\times$};
  \node[class1, font=\tiny] at (0.55,0.15) {$\times$};
  \node[class1, font=\tiny] at (1.05,0.25) {$\times$};
  \node[class1, font=\tiny] at (0.35,0.32) {$\times$};
  \node[class1, font=\tiny] at (0.85,0.2)  {$\times$};
  \node[class1, font=\tiny] at (1.3,0.35)  {$\times$};
  \draw[testred, thick] plot[smooth] coordinates {(-0.1,0.65) (0.2,0.82) (0.45,0.52) (0.7,0.78) (0.95,0.55) (1.2,0.72) (1.5,0.62)};
}}

% One "hypothesis space histogram" panel (used twice on the Meta Level
% slide: once for the prior $\Phi(\varphi)$, once for the posterior
% $\Phi(\varphi \mid D_n)$) -- a bar over $\varphi$ per representative
% hypothesis, with its decision-boundary icon underneath.
%
% Usage: \metapanel{<y-axis label>}{<bar1 height>}{<bar2>}{<bar3>}{<bar4>}
%   e.g. \metapanel{$\Phi(\varphi)$}{0.9}{1.0}{0.85}{0.8}
%
% HOW TO ADJUST:
%   - #1 = the y-axis label (put $\Phi(\varphi)$ or $\Phi(\varphi\mid
%     D_n)$ here when calling it).
%   - #2..#5 = the four bar heights, in that order; these are the ONLY
%     thing that differs between the "prior" and "posterior" calls on the
%     Meta Level slide -- raise #3 (hypLinB's bar), say, to show the
%     posterior concentrating on that hypothesis.
%   - Each bar/icon pair is one hardcoded block: "rectangle (x0,0)
%     rectangle (x1,height)" draws the bar, "\node at (xc,-0.85)
%     {\hyp...}" places its icon below the axis. To add a 5th hypothesis,
%     copy one block, shift its x-coordinates past 8.65, extend the axis
%     \draw's endpoint (currently 8.8), and add a 6th macro argument.
%   - "$\cdots$" nodes are just the ellipsis marks between bars,
%     indicating "more hypotheses in between" -- move them if you resize
%     the gaps.
%   - "scale=0.5" shrinks the whole panel to fit two of them side by side
%     on one slide; raise/lower it (together with the minipage height and
%     column widths in main.tex) if you change how many panels you show.
\newcommand{\metapanel}[5]{%
  \begin{tikzpicture}[scale=0.5]
    \draw[-{Latex}] (-1.9,0) -- (10,0) node[right] {$\varphi$};
    \draw[-{Latex}] (-1.9,-0.2) -- (-1.9,3.7) node[above] {#1};
    
    \draw[fill=testred!70, draw=testred] (-1.2,0) rectangle (0.2,#2);
    \node at (-0.5,-1.85) {\scalebox{0.8}{\hypLinA}};
    \node at (1.0,0.35) {$\cdots$};
    
    \draw[fill=testred!70, draw=testred] (1.8,0) rectangle (3.2,#3);
    \node at (2.5,-1.85) {\scalebox{0.8}{\hypLinB}};
    \node at (4.05,0.35) {$\cdots$};
    
    \draw[fill=testred!70, draw=testred] (4.8,0) rectangle (6.2,#4);
    \node at (5.5,-1.85) {\scalebox{0.8}{\hypCurve}};
    \node at (7.1,0.35) {$\cdots$};
    
    \draw[fill=testred!70, draw=testred] (7.8,0) rectangle (9.2,#5);
    \node at (8.5,-1.85) {\scalebox{0.8}{\hypWiggly}};
  \end{tikzpicture}%
}
